\section{統計の誤用とデータの倫理}

統計は, 複雑な現象から有意味な構造を抽出するための道具であるが, その力は両刃の剣である.  
適切に使えば意思決定を支える礎となる一方, 使い方を誤れば現実を歪める.  

人間はしばしば, 自分の直感や先入観を裏付ける情報だけを強く記憶し, 不都合な情報を軽く扱う認知バイアスを持つ.  
さらに組織では, 短期的成果や都合のよい物語が選ばれやすい構造的インセンティブが働く.  
この二重構造こそが, 統計の誤用を生む温床である.

データドリブンが当たり前になった社会では, 統計の誤用が組織の判断を長期的に損ね, 社会的信頼や競争力を削ぐ事例が後を絶たない.  

本章では, 図表の設計や統計的推論の現場で陥りやすい典型的な誤りを数値と事例で点検し, その背後にある心理的・構造的要因と, 実務での予防策を明確にする.  
これらの対策は, 単に誤りを防ぐだけでなく, 隠れていた価値ある情報を発見し, 公正で持続的な意思決定を支える力にもなる.

\subsection{図表の設計がもたらす誤解}

図やグラフは, 数値の比較や変化を直感的に理解させる強い効果を持つ.  
しかし, 軸やスケール, 表示期間といった仕様の選び方ひとつで印象は大きく変わる.  
これは, 人間の視覚が絶対値ではなく相対的な差や傾きを重視する性質による.

ここでは, 「なぜ誤解が生じるのか」を構造的に示し, 読み手・作り手の双方に必要な確認基準を整理する.

\subsubsection{軸・スケール・期間の設定}

縦軸の起点が0でない棒グラフは, 棒の長さの比が数値の比と一致せず, 差を過大に見せる.  
また, 目盛の間隔を不均一にしたり, 都合の良い短期間だけを切り取ることで, 勾配や変化の印象を操作できてしまう.  

例えば, 気温が$+0.5^\circ$上昇した場合, 縦軸を0--5℃で描くのと0--1℃に圧縮するのとでは, 視覚的な差の比は何倍も異なる.  
同じデータで縦軸範囲を変えた2種類のグラフを並べて見たとき, その印象がどのように変わるかを想像してみると, この効果を直感的に理解できる.

\subsubsection{部分的なデータ提示}

チェリーピッキングとは, 技術的には正しい一部のデータだけを切り出し, 全体像を歪める行為である.  
成功例のみの提示や, 「統計的に有意でない」を理由に不都合な結果を無視することが典型である.  

生存者バイアスはその極端な例で, 帰還機の被弾箇所だけを見て補強点を誤る話は有名だ.  
これはあたかも, 雨の日に傘を持っていた人だけを見て「みんな濡れない」と結論づけるようなものである.  
全体の情報を示す場合と部分だけを切り出す場合で, 読み手の判断がどのように変わるかを想像してみると, 情報選択の影響がより鮮明になる.

\subsubsection{図種や装飾による印象操作}

面積比較が難しい円グラフは比率の差を直感的に誤認させやすい.  
3D化された棒や円は遠近感で面積が歪み, 定量的判断を阻害する.  
二軸グラフはスケール設定次第で相関があるように見せかけられる.  
配色は無意識の評価を誘導するため, 色選びにも慎重さが要る.

\subsubsection{読み手が確認すべきこと}

出所は明示されているか.  
縦軸の起点と目盛間隔は適切か.  
単位や期間に選択理由が説明されているか.  
図種は目的（比較・推移・構成比）と合致しているか.  
欠落しているデータは何か.  

こうした点を一つ一つ確認した場合としなかった場合で, 誤解の発生可能性がどの程度変わるかを想像すると, その重要性がより明確になる.

\subsubsection{作り手が必ず開示すべき条件}

収集過程, 前処理手順, サンプルの代表性, 欠測や外れ値の扱い, 仕様変更履歴を明記し, 再現性を担保する.  
見た目の洗練はその後の話であり, 条件と限界の明示が先である.

\subsection{サンプリングの偏りと生存者バイアス}

標本調査の前提は, 標本が母集団を代表することである.  
無作為抽出を外れると, 推定は系統的に歪む.  

典型例として, 時間帯や場所が偏った街頭調査, 固定電話所有に偏りがある時代の電話調査, ネットの自己選択アンケートがある.  
成果が良好な層にのみ依頼する満足度調査は, それ自体が偏りを内包する.  

\textbf{用語補足}: MCAR（完全無作為欠測）, MAR（条件付き無作為欠測）, MNAR（無作為でない欠測）は, 欠測データの発生メカニズムを分類する用語である.

生存者バイアスは, 選抜を生き残った対象だけから結論を出す誤謬である.  
戦時の帰還機の例, 成功企業のベストプラクティスだけを収集する事例, 会員継続者だけの満足度調査などが該当する.  
調査対象を「生き残り」に限定した場合とそうでない場合で, 結果や結論がどう変わるかを想像してみると, この偏りの影響がより実感できる.

\subsection{統計的推論に潜む罠}

これまでは, データの取得や表示の段階で生じる歪みを見てきた.  
しかし, たとえデータが正しく集められ, 見た目にも公平に表示されていたとしても, 推論の段階で誤解が生じれば結論は簡単に誤る.  
ここからは, その「推論段階の落とし穴」に踏み込む.

\subsubsection{シンプソンのパラドックス}

各層で同方向の差が, 集約後に逆転する.  
原因は交絡因子の分布差である.  

もし集約結果と層別結果の両方を比較できたとしたら, どのように結論が変わるかを考えてみると, この現象の重要性が理解しやすい.

\subsubsection{P値の誤解と再現性危機}

P値は「帰無仮説の下で観測以上に極端な結果が出る確率」であり, 効果の大きさでも仮説の真偽の確率でもない.  
大標本では実務的に些末な差でも有意になりやすい.  

出版バイアスと$p$ハッキング（都合の良い結果が出るまで分析を繰り返す行為）は再現性危機を助長する.  
これは, まるで料理の味見を何十回もして, たまたま美味しい一口が出たときだけをレシピとして記録するようなものである.  

仮に同じデータで異なる分析者が異なる手法を選んだ場合, 得られる結論がどれほど変わり得るかを想像することで, この問題の深刻さがより明確になる.

\subsubsection{平均値への過信}

平均は外れ値に敏感で, 歪んだ分布では典型値を表さない.  
所得, 都市規模, 企業規模のように重い尾を持つ分布では, 中央値や分位点, ロバストな散布指標を併記する.  

平均値だけで判断した場合と, 中央値などを併記した場合で意思決定がどう変わるかを想像すると, この指標選択の重要性が理解しやすい.

\subsubsection{多重比較とデータ漁り}

独立な検定を$m$回行うと, 少なくとも1回偽陽性が出る確率は$1-(1-\alpha)^m$である.  
例えば$\alpha=0.05, m=20$で約0.64となる.  

もし無数の仮説検定を繰り返した場合, どの程度「偶然の発見」が増えるのかを想像してみると, この問題が直感的に理解できる.

\subsubsection{相関と因果の混同}

相関は因果を含意しない.  
交絡因子や逆因果が疑似相関を生む.  

気温上昇はアイスの売上と水難事故を同時に増やすが, 因果はない.  
因果関係がある場合とない場合で, 施策や判断がどう変わるかを考えてみると, 因果識別の重要性がわかる.

\subsection{公的統計の信頼性から得る教訓}

ここまで扱ったのは主に個人や組織内部での誤用であった.  
しかし, 統計の誤用が制度的レベルで発生すると, その影響は桁違いに大きく, 社会全体に及ぶ.  
その典型が, 公的統計の不適切運用である.

毎月勤労統計問題では, 本来全数調査すべき大規模事業所に抽出調査を適用し, さらに復元処理を怠った.  
結果として平均賃金が過小に算出され, 延べ約2{,}000万人に総額500億円超の過少給付が生じた.  

教訓は明確である.  
調査設計の逸脱は系統的誤差を生み, 社会保障の配分という実体的結果に直結する.  
手続の遵守, 設計変更の記録と公開, 外部検証の受容が不可欠である.  

同じ制度を, 設計通り正確に運用した場合と, 手順を省略・変更した場合とで, 最終的な給付額や対象者の分布がどのように変わるかを想像すると, この種の逸脱がもたらす影響の大きさが実感できる.

\subsection{倫理原則とその運用}

\subsubsection{三本柱}

倫理は, \textbf{プライバシーの保護}, \textbf{公平性}, \textbf{透明性と説明責任}に集約される.  

個人情報の取り扱いでは, 目的の特定, 同意, 安全管理(暗号化, 権限管理), 匿名化/仮名化が基本である.  
公平性では, 学習データの代表性確保, バイアス検知・緩和, 多様なステークホルダーによるレビューを標準化する.  
透明性と説明責任では, 判断の根拠を説明可能にし, 役割分担と責任の所在を事前に定義する.  

これらの原則を満たす場合と満たさない場合で, 組織や社会の信頼度がどのように変化するかを想像すると, その重要性が具体的に見えてくる.

\subsubsection{運用枠組み}

FAT原則(Fairness, Accountability, Transparency)を組織運用に落とし込む.  
具体的には, 説明可能なAI(XAI: eXplainable AI)の導入, 倫理委員会と外部監査の設置, 仕様変更と意思決定の監査記録, インシデント対応計画の整備である.  

研究では, オープンサイエンス(データ, コード, 手順の公開)と事前登録を基本とし, 事後探索は探索と明示する.  
人を対象とする場合は, インフォームド・コンセントの理念を守り, 撤回可能性を保障する.  

もしこうした運用枠組みが存在する組織と存在しない組織を比較したら, 長期的な成果や外部からの評価にどのような差が生まれるかを想像すると, 整備の必要性がより鮮明になる.

\subsection{現場で使う点検表}

\begin{itemize}
  \item \textbf{データの出所}. 取得経路, ライセンス, 目的を明記しているか.
  \item \textbf{サンプリング}. 抽出方法, 回答率, 欠測の発生機構(MCAR/MAR/MNAR)を記述しているか.
  \item \textbf{前処理}. 外れ値基準, 欠測処理, 集計単位変更の手順が再現可能か.
  \item \textbf{指標}. 平均に加えて中央値, 分位点, 散布, 効果量, 信頼区間を併記しているか.
  \item \textbf{検証}. 多重比較補正, 交差検証, 感度分析(仮定変更)を実施しているか.
  \item \textbf{図表}. 軸起点と目盛, 単位, 期間, 図種の適合性を点検し, 代替表示を試したか.
  \item \textbf{倫理}. 目的外利用の防止, 権限最小化, ログと監査, 連絡窓口の明示があるか.
\end{itemize}

この点検表の各項目を満たした場合と満たさなかった場合のアウトプットや判断の質の違いを想像すると, 日常的な確認作業がいかに実務の信頼性を高めるかが理解しやすい.

\subsection{まとめ}

本章は, 図表設計, サンプリング, 推論, 公的統計の運用, 倫理原則という5つの層で, 誤用の構造と対策を定量的に確認した.  

重要なのは, 誤用の回避を防衛的課題としてではなく, 不確実な世界でより良い意思決定を実現するための積極的基盤として位置付けることである.  

技術的厳密さと倫理的配慮を同時に満たすことが, 信頼できる統計実務の必要条件であり, さらにそれは, 新しい価値や機会を生み出す源泉にもなり得る.  

もしこの章で挙げた対策を一貫して実践できたなら, 組織や社会にどのような変化が起こり得るかを思い描くことは, 学びを実践へとつなぐ第一歩となる.
